\section{Nomenclature}
\label{sec:nomenclature}

\subsection{Notation Policy}

This notebook uses a strict notation policy consistent with the C++ module:
\begin{itemize}
  \item scalar variables use italic math symbols, for example \(P_i, Q_i, V_i\),
  \item sets use calligraphic symbols, for example \(\mathcal{N}_{ac}, \mathcal{E}_{dc}\),
  \item all semantic qualifiers use upright superscripts via \(\mathrm{\cdot}\), for example \(P_i^{\qual{spec}}, P_i^{\qual{calc}}, V_k^{\qual{dc}}\),
  \item for variable symbols, subscripts are reserved for ordered indices (bus/component/time) and not for semantic class labels,
  \item set symbols may use compact domain tags in subscript, e.g., \(\mathcal{N}_{ac}\), \(\mathcal{E}_{dc}\),
  \item model-family and semantic words must stay in superscripts: write \(P^{\qual{dcdc}}\), \(P^{\qual{flex}}\), \(P_r^{\qual{rated}}\), \(c_r^{\qual{curt}}\), and \(\mathcal{R}^{\qual{curt}}\).
\end{itemize}

\subsection{Index and Qualifier Conventions}

\begin{longtable}{L{0.24\linewidth}L{0.66\linewidth}}
\toprule
Token & Meaning \\
\midrule
\endhead
\(i,j\) & AC bus indices \\
\(k,m\) & DC bus indices \\
\(c\) & converter index (VSC / DCDC) \\
\(b\) & branch index \\
\(g\) & generator index \\
\(\qual{spec}\) & specified (assembled) quantity from component injections/demands \\
\(\qual{calc}\) & network-calculated quantity from current state \\
\(\qual{ac}\), \(\qual{dc}\) & AC-side / DC-side quantity qualifier \\
\(\qual{conv}\) & converter-side quantity qualifier \\
\(\qual{dcdc}\), \(\qual{flex}\) & DCDC-dispatch / flexible-load variable class qualifier \\
\(\qual{rated}\), \(\qual{curt}\) & rated-capacity / curtailment-related qualifier \\
\(\qual{min}\), \(\qual{max}\) & lower/upper operating bounds \\
\bottomrule
\end{longtable}

\subsection{Canonical Naming Grammar}

Use the canonical symbol template
\begin{align}
X_{\text{indices}}^{\text{qualifiers}},
\end{align}
with the following rules:
\begin{itemize}
  \item index block (\(_{\cdot}\)): ordered indices only (\(i,j,k,c,r,t,\dots\)),
  \item qualifier block (\(^{\cdot}\)): semantic tags only (\(\qual{spec}\), \(\qual{calc}\), \(\qual{conv}\), \(\qual{dcdc}\), \(\qual{flex}\), \(\qual{rated}\), \(\qual{curt}\), \(\qual{min}\), \(\qual{max}\), \dots),
  \item when multiple qualifiers are needed, order them as: model/class \(\rightarrow\) side or direction \(\rightarrow\) status/bound/cost.
\end{itemize}

Examples:
\begin{align}
P_i^{\qual{conv,ac}},\quad
P_k^{\qual{dcdc,in}},\quad
P_r^{\qual{rated}},\quad
c_r^{\qual{curt}}.
\end{align}

\subsection{Superscript/Subscript Rule}

\begin{longtable}{L{0.30\linewidth}L{0.60\linewidth}}
\toprule
Rule & Usage \\
\midrule
\endhead
Subscript (\(_{\cdot}\)) & for variable symbols: ordered indices only (\(i,j,k,c,r,t\)); set symbols may use domain tags (\(\mathcal{N}_{ac}\), \(\mathcal{E}_{dc}\)). \\
Superscript (\(^{\cdot}\)) & semantic labels: domain, role, operating mode, variable class, bounds, and status qualifiers. \\
Forbidden mixed style & avoid semantic words in subscript, such as \(P_{dcdc}\), \(P_{flex}\), \(\mathcal{R}_{curt}\), \(P_{rated}\), \(c_{curt}\). \\
Preferred style & \(P^{\qual{dcdc}}, P^{\qual{flex}}, \mathcal{R}^{\qual{curt}}, P_r^{\qual{rated}}, c_r^{\qual{curt}}\). \\
\bottomrule
\end{longtable}

\subsection{Integration Status Codes}

The integration matrix uses these codes:
\begin{itemize}
  \item \statuscode{D}: direct model consumption in the module,
  \item \statuscode{F}: fixed-only consumption (no explicit variable),
  \item \statuscode{P}: consumed after projection to canonical equivalents,
  \item \statuscode{N}: not consumed by the module.
\end{itemize}
